\documentclass[10pt,a4paper]{amsart}
\usepackage[margin=19mm]{geometry}
\usepackage{amsmath,amssymb,mathtools}
\usepackage[scr=rsfs,cal=euler]{mathalfa}
\usepackage{xfrac}
\usepackage[unicode,hidelinks,bookmarks=false]{hyperref}
\newcommand{\ie}{i.e.\ }
\newcommand{\cf}{cf.\ }
\newcommand{\resp}{respectively\ }
\newcommand{\Subsection}{\subsection}
\providecommand{\nobreakdash}{\nobreak\hbox{-}}
\setlength{\emergencystretch}{2em}
\multlinegap0pt
\usepackage{bm}
\usepackage{bbm}
\usepackage{dsfont}
\usepackage{xspace}
\usepackage{cancel}
\usepackage{mathtools}
\usepackage{amsthm}
\makeatletter
\@ifundefined{proposition}{\newtheorem{proposition}{Proposition}}{}
\makeatother
\title[Nonlinear Nonlocal Diffusion Equations for the Analysis of C]{Nonlinear Nonlocal Diffusion Equations for the Analysis of Continuous Coordination and Anti-Coordination Type Games}
\author{John S. McAlister, Nina H. Fefferman, Tadele A. Mengesha}
\date{Source: arXiv:2506.13929v1 (CC BY 4.0)}
\begin{document}
\maketitle

\noindent\emph{Source and licence.} Nonlinear Nonlocal Diffusion Equations for the Analysis of Continuous Coordination and Anti-Coordination Type Games. Version arXiv:2506.13929v1 (CC BY 4.0), distributed under the Creative Commons Attribution 4.0 International licence, as recorded in the arXiv licence metadata for this exact version and shown on the abstract page https://arxiv.org/abs/2506.13929v1. Full rights notice: licenses/arxiv-2506.13929-CC-BY-4.0.txt.\par\smallskip

\noindent\textit{Editorial scope.} Counted unit: the Proposition whose source label is \texttt{PROP:ModelWellFounded} in the article (source0.tex lines 182--192, source label \texttt{PROP:ModelWellFounded}), delivered together with the article's own complete original proof of it (source0.tex lines 194--271, one \texttt{proof} environment of 78 lines). No mathematics is written, repaired or re-derived: statement and proof are the article's own text line for line. Required context retained uncounted: ctsStratPayoff (lines 166--168). Every definition, display and label of those lines is retained unchanged. Omitted from this package and not redistributed: 1--165, 169--181, 193--193, 272--1264. Nothing inside a retained excerpt is omitted or altered: every retained body is delivered exactly as the article's lines are written, so no omission receipt of any kind accompanies this package. Citations: the retained excerpts carry no citation command at all, so no bibliography entry of the article is transcribed and no reference is redistributed. Reference adaptation: none was needed and none was made; every reference inside the delivered unit resolves inside it, so no unresolved reference or citation is produced. Printed slips kept literal: no printed word, symbol, hypothesis or number of the counted statement or of its proof was altered. Layout and class: the article's own class is \texttt{revtex4-2} with packages including graphicx, esint, tikz, amsthm, amssymb, hyperref, ulem; the wrapper is the lane's pinned \texttt{amsart} editorial wrapper at 10pt on A4, which restates the theorem-like environments the retained text needs and adds the article's own macro definitions that the retained text needs (none), with extra wrapper packages (bm, bbm, dsfont, xspace, cancel, mathtools). The delivered numbering is the unit's own. Assets: no retained span contains a figure environment, a graphics-inclusion command, a PDF-page inclusion, a table or a TikZ picture; no image file is copied into this package, the package declares no asset dependency and no image file is redistributed with it. Licence: version arXiv:2506.13929v1 of this article is distributed under the Creative Commons Attribution 4.0 International licence, as recorded in the arXiv licence metadata for this exact version and shown on the abstract page https://arxiv.org/abs/2506.13929v1. A full rights notice is delivered at licenses/arxiv-2506.13929-CC-BY-4.0.txt. Characters: every retained excerpt is pure ASCII, so the delivered engine, XeLaTeX with its default Latin Modern text font, reports no missing character.
\begin{equation}\label{ctsStratPayoff}
w(x|u)=\int_{\Omega}K(x,y)\rho(u(x)-u(y))dy.
\end{equation}

    \begin{proposition}\label{PROP:ModelWellFounded}
        Under myopic best response, bounded strategy profiles of the game with players $\Omega\subseteq \mathbb{R}^n$ choosing strategies in $\mathbb{R}$, with fitness as in \eqref{ctsStratPayoff}, will evolve as
       
       \[\frac{\partial}{\partial t}u(x,t)=\int_{\Omega}K(x,y)\rho'(u(x,t)-u(y,t))dy\]
        so long as the following three hypotheses are met
        \begin{itemize}
            \item[(H1)] Players change their strategies in arbitrarily small time steps $\Delta t$
            \item[(H2)] Players incur a quadratic cost $\frac{h^2}{\Delta t}$ for changing their strategy
            \item[(H3)] $\rho\in C^{1,1}(\mathbb{R})$ and $\rho(z)\leq Cz^2+A$ for some nonnegative $C,A$
        \end{itemize}
    \end{proposition}

    \begin{proof}
        Consider a bounded strategy profile $u(\cdot, t):\Omega\to \mathbb{R}$. Every player will seek to update their strategy by some amount $h$ in order to take on their best response to $u(\cdot, t)$ after a time step of $\Delta t$ and in doing so they will incur a cost of $\frac{h^2}{2}$. Let 
        
       \[S_x(h):=\int_\Omega K(x,y)\rho(u(x,t)+h-u(y,t))dy\] be the payoff that player $x$ will receive after updating their strategy by $h$. Because $\rho$ is subquadratic (H3), we can see that 

        \begin{equation*}
        \begin{split}
            S_x(h)&\leq C\int_\Omega K(x,y)(h+(u(x,t)-u(y,t))^2dy+\underbrace{A\sup_{x\in\Omega}\|K(x,\cdot)\|_{L^1(\Omega)}}_{A_1}\\
            &\leq \underbrace{C\int_\Omega K(x,y)dy}_{C_1}h^2 + 2hC\int_\Omega K(x,y)(u(x,t)-u(y,t))dy\\
            &\quad +C\int_\Omega K(x,y)(u(x,t)-u(y,t))^2dy+A_1
        \end{split}
        \end{equation*}
        Since $K\in C_b^0(\Omega; L^1(\mathbb{R}^n))$, $C_1$ and $A_1$ can be bounded by constant independent of $x$. Now because $u$ is bounded we know that $|u(x,t)-u(y,t)|$ is bounded by some $M$ and thus we can write 
        \begin{equation*}
            \begin{split}
                S_x(h)&\leq C_1h^2+\underbrace{2CM\sup_{x\in\Omega}\|K(x,\cdot)\|_{L^1(\Omega)}}_{B_1}h+ \underbrace{CM^2\sup_{x\in\Omega}\|K(x,\cdot )\|_{L^1(\Omega)} + A_1}_{A_2}\\
                &\leq C_1h^2+B_1h+A_2
            \end{split}
        \end{equation*}
        Having shown that $S_x(h)$ is uniformly subquadratic in $h$, we also note that when $\rho\in C^{1,1}$ (H3), $S_x(h)$ is also in $C^{1,1}$. Observe that 
        \begin{equation*}
            \frac{d}{dh}S_x(h)=\int_\Omega K(x,y)\rho'(u(x,t)+h-u(y,t))dy
        \end{equation*}
        And so clearly for any compact subdomain $I\subset \mathbb{R}$ we can write 
        \begin{equation}\label{EQ:SinC11}
        \begin{split}
            \bigg|&\frac{d}{dh}S_x(h_1)-\frac{d}{dh}S_x(h_2)\bigg|\\
            &\leq \int_\Omega |K(x,y)| |\rho'(u(x,t)+h_1-u(y,t))-\rho'(u(x,t)+h_2-u(y,t))|dy\\
            &\leq \int_\Omega K(x,y)L_{\rho}|h_1-h_2|dy\\
            &\leq \sup_{x\in\Omega}\|K(x,\cdot)\|_{L^1(\Omega)}L_{\rho}|h_1-h_2|
            \end{split} 
        \end{equation}
        Because $u$ is bounded, the choice of a compact subdomain $I$ gives us a Lipschitz constant for $\rho'$ on the subdomain $[-2\sup |u|-\sup I,2\sup|u|+\sup I]$ which is called $L_\rho$ in \eqref{EQ:SinC11}. Therefore $\frac{d}{dh}S_x(h)$ is locally Lipschitz. 

        This is important because it means that, for any $x$, when $\Delta t$ is small enough, $S_x(h)-\frac{h^2}{\Delta t}\leq C_2h^2+A_3$ for some negative $C_2$ and some $A_3$. Thus $S_x(h)-\frac{h^2}{\Delta t}$ must have a global maximizer, $h^*$ and that global maximizer will satisfy 
        
       \[\frac{d}{dh}S_x(h^*)=2\frac{h^*}{\Delta t}\]
        because everything is continuously differentiable. 

        Having observed this, we note also that there must be a negative $h^-$ so that when $h<h^-$ then $S_x(h)-\frac{(h)^2}{\Delta t}<S_x(0)$. Likewise there must be an $h^+$ so that $h>h^+\implies S_x(h)-\frac{(h)^2}{\Delta t}<S_x(0)$.
        These $h^-$ and $h^+$ will form a compact interval which will surely contain $h^*$. Moreover, from the chosen interval we have a Lipschitz constant for $\frac{d}{dh}S_x$ which we call $L_S$.
        The next step is to put bounds on $h^*$, in order to do this we will consider different cases: when $h^*>0, h^*<0$, and $h^*=0$. 

        If $h^*>0$ we write that 
        \begin{equation*}
            \begin{split}
                -h^*L_S&\leq \frac{d}{dh}S_x(h^*)-\frac{d}{dh}S_x(0)\leq h^*L_S\\
                -h^*L_S&\leq \frac{h^*}{2}-\frac{d}{dh}S_x(0)\leq h^*L_S\\
                \end{split}
        \end{equation*} 
        Reorganization on each side of the inequality, with the knowledge that $\Delta t$ can be made small enough so that $2-L_S\Delta t>0$, will yield 
        \begin{equation*}
                \frac{d}{dh}S_x(0)\frac{\Delta t}{2+L_S\Delta t}\leq h^*\leq \frac{d}{dh}S_x(0)\frac{\Delta t}{2-L_S\Delta t}
        \end{equation*}

        Likewise we can show that if $h^*<0$ we will get the inequality 
        \begin{equation*}
            \frac{d}{dh}S_x(0)\frac{\Delta t}{2-L_S\Delta t}\leq h^*\leq \frac{d}{dh}S_x(0)\frac{\Delta t}{2+L_S\Delta t}
        \end{equation*}

        We also note here that if $h^*=0$ this requires that $\frac{d}{dh}S_x(0)=0$. 

        Now that we have bounds on $h^*$, recall that $h^*$ is the change in strategy in one time step. That it $h^*=u(x,t+\Delta t)-u(x,t)$. If we make this substitution and divide our inequalities by $\Delta t$ we see that 
        \begin{equation*}
            \begin{split}
                \frac{d}{dh}S_x(0)\frac{1}{2+L_S\Delta t}&\leq \frac{u(x,t+\Delta t)-u(x,t)}{\Delta t}\leq\frac{d}{dh}S_x(0)\frac{1}{2-L_S\Delta t}\quad \quad h^*>0\\
                \frac{d}{dh}S_x(0)\frac{1}{2-L_S\Delta t}&\leq \frac{u(x,t+\Delta t)-u(x,t)}{\Delta t}\leq\frac{d}{dh}S_x(0)\frac{1}{2+L_S\Delta t}\quad \quad h^*<0
            \end{split}
        \end{equation*}
        Both inequalities trivially hold in the case that $h^*=0$. In any case, when we take $\Delta t\to 0$ we see that by squeeze theorem 
        \begin{equation*}
            \frac{\partial}{\partial t}u(x,t)=\frac{1}{2}\frac{d}{dh}S_x(0)
        \end{equation*}

        To complete the proof we need only note that 
        
       \[\frac{d}{dh}S_x(0)=\int_\Omega K(x,y)\rho'(u(x,t)-u(y,t))dy\] and do a trivial rescaling of space-time to arrive at the desired nonlocal equation.
        \end{proof}

\end{document}
